\section{The genuine locally finite divisor sum}
\label{sec:hironaka-locally-finite-divisor-sum}

Let $D$ be an actual effective divisor supported on an SNC boundary $E$.
The members of $E$ may be disconnected, and the ambient manifold need
not be compact.  Multiplicities belong to the genuine connected carriers
$\mathcal C(E)$, rather than to the original list members.  Their local
finiteness and the intrinsic monomial identity give an actual divisor
sum on a neighbourhood of every point.

\begin{lemma}[Powers depend on the whole stalk]
\label{hir:ideal-power-stalk-congruence}
\lean{Hironaka.HironakaLocallyFiniteDivisorSum.idealPow_mem_iff_of_mem_iff}
\leanok
Suppose $I_x=J_x$ as complete holomorphic ideal stalks.  Then
$(I^r)_x=(J^r)_x$ for every $r\ge0$.  In particular, if two closed sets
agree near $x$, all powers of their actual vanishing ideals agree at $x$.
\end{lemma}
\begin{proof}
For $r=0$, both powers are the holomorphic unit ideal.  If the assertion
holds for $r$, a generator of $(I^{r+1})_x$ is a product of a member of
$(I^r)_x$ with a member of $I_x$.  Both factors belong to their corresponding
$J$ stalks, so the product belongs to the generating family for
$(J^{r+1})_x$.  This also holds for every finite holomorphic linear
combination, and hence gives inclusion of the whole product stalk.
Interchanging $I$ and $J$ gives equality.  For vanishing ideals, equality
of the sets near $x$ makes the eventual vanishing condition identical
for every holomorphic germ, giving the required original stalk equality.
\end{proof}

\begin{lemma}[Support of an actual finite hypersurface sum]
\label{hir:finite-component-sum-support-containment}
\uses{hir:component-product-effective-sum,hir:effective-divisor-multiples}
\lean{Hironaka.HironakaLocallyFiniteDivisorSum.support_componentSum_subset}
\leanok
For a finite weighted family $A=((C_i,b_i))$ of closed smooth
hypersurfaces, the actual effective sum $G=\sum_i b_i C_i$ satisfies
$\operatorname{Supp}G\subseteq\bigcup_i C_i$ and
$I_G=\prod_i I_{C_i}^{b_i}$.
\end{lemma}
\begin{proof}
The genuine hypersurface divisor has ideal $I_{C_i}$ and support $C_i$.
A zero multiple is the zero divisor and has empty support; a positive
multiple retains support $C_i$.  Actual divisor addition is multiplication
of invertible ideals.  Its support is the union of the two supports:
at a point outside both zero loci, each coherent factor contains a unit
germ, whose product is a unit; at a point in either zero locus, every
product germ vanishes.  Induction on the finite list proves the support
containment.  The same induction, using the actual power identity for
natural multiples, proves the stated complete ideal identity.
\end{proof}

\begin{lemma}[One actual coefficient on every genuine component]
\label{hir:locally-finite-original-divisor-coefficients}
\uses{hir:connected-boundary-component-multiplicities,hir:connected-boundary-components-locally-finite}
\lean{Hironaka.HironakaLocallyFiniteDivisorSum.exists_original_coefficients_boundaryComponents}
\leanok
In positive ambient dimension there is a single function
$a:\mathcal C(E)\to\mathbb N$ such that, for every $C\in\mathcal C(E)$
and every original point $y\in C$,
$e_C(I_D,y)=a(C)$.  The family of the actual carriers is locally finite.
\end{lemma}
\begin{proof}
The actual connected-component exponent theorem supplies $a(C)$ and the
equivalence $\operatorname{IsDivNear}(I_D,C,r,y)\Longleftrightarrow r\le a(C)$
at every point of $C$.  Invertibility gives a nonzero germ in each stalk
of $I_D$, so the genuine finite divisor exponent characterizes exactly
the same divisibility thresholds.  Taking $r=a(C)$ and
$r=e_C(I_D,y)$ gives the two inequalities and therefore equality.
The local finiteness theorem applies to the genuine connected boundary
carriers without any compactness assumption.
\end{proof}

\begin{lemma}[The complete stalk dictionary for a finite local refinement]
\label{hir:monomial-local-component-product-dictionary}
\uses{hir:ideal-power-stalk-congruence,hir:finite-component-sum-support-containment,hir:component-product-original-exponent,hir:intrinsic-supported-ideal-monomial,hir:boundary-component-germ-divisibility}
\lean{Hironaka.HironakaLocallyFiniteDivisorSum.monomialPart_mem_iff_componentPowerProduct_of_refinement}
\leanok
Let $F$ be a finite list of distinct genuine components of $E$ containing
every component through $x$.  Assign to $C\in F$ the actual original
coefficient $a(C)$.  Then the entire intrinsic monomial stalk of $I_D$
at $x$ equals the actual product stalk $\prod_{C\in F}I_C^{a(C)}$.
\end{lemma}
\begin{proof}
Every original member $H_i$ through $x$ has a genuine component $C$
through $x$.  It belongs to $F$.  The two carriers agree on a neighbourhood
of $x$, so their actual vanishing-ideal powers have the same whole stalk.
Their genuine divisor exponents are equal by the component-to-parent
germ bridge, and are $a(C)$.  Conversely, every $C\in F$ through $x$
has an original parent $H_i$ with the same germ at $x$.  Thus the
definition of the intrinsic monomial stalk, as simultaneous membership
in all the relevant powered vanishing ideals, is unchanged by replacing
$E$ with $F$.

Form the actual effective sum $G=\sum_{C\in F}a(C)C$.  The preceding
support lemma makes it SNC-supported, and its full ideal is the actual
component-power product.  The genuine weighted SNC exponent calculation
identifies $e_C(I_G,x)=a(C)$ for every $C$ through $x$, including crossing
points and zero coefficients.  Hence the intrinsic monomial stalks of
$I_D$ and $I_G$ relative to $F$ are identical.  The actual SNC-supported
divisor identity gives $I_G=\mathcal M(I_G,F)$.  Combining these complete
stalk identities proves the claim.  No ideal factorization of $D$ was
assumed in this argument.
\end{proof}

\begin{theorem}[Positive local finite products with original coefficients]
\label{hir:locally-finite-original-component-product}
\uses{hir:locally-finite-original-divisor-coefficients,hir:monomial-local-component-product-dictionary,hir:intrinsic-supported-ideal-monomial}
\lean{Hironaka.HironakaLocallyFiniteDivisorSum.exists_positive_locallyFinite_original_component_products,Hironaka.HironakaLocallyFiniteDivisorSum.exists_locallyFinite_original_component_products}
\leanok
In every ambient dimension, there is one actual coefficient function
$a:\mathcal C(E)\to\mathbb N$ with
$a(C)=e_C(I_D,y)$ for all $y\in C$.  The carrier family is locally finite.
For every $x$ there are an open neighbourhood $U$ and a finite list of
distinct actual carriers with positive coefficients,
$A=((C_i,a(C_i)))$, such that for every $y\in U$ and every holomorphic
germ represented by $g$,
\[
 g\in(I_D)_y\quad\Longleftrightarrow\quad
 g\in\left(\prod_i I_{C_i}^{a(C_i)}\right)_y.
\]
\end{theorem}
\begin{proof}
First suppose the ambient dimension is positive.  Select the single
genuine function $a$ above.  Local finiteness gives an open neighbourhood
$U$ of $x$ meeting only finitely many genuine boundary components.
Enumerate those components once to obtain $F$.  Their genuine geometry
and the finite-refinement theorem make $F$ SNC.  At any $y\in U$, every
genuine component containing $y$ meets $U$ and therefore belongs to $F$.
The preceding dictionary applies at every such $y$.  Since $D$ is
actually SNC-supported, $I_D=\mathcal M(I_D,E)$, which proves the desired
whole-stalk equality with the full weighted product on $F$.

Delete exactly those pairs whose coefficient is zero.  Every deleted
factor is $I_C^0=\mathcal O$.  The remaining products are coherent, so
multiplication by the holomorphic unit ideal leaves the entire actual
product ideal unchanged.  Induction on the weighted list verifies this
without changing any stalk.  The retained coefficients are positive,
remain values of the same global function $a$, and the retained carriers
remain distinct.

In dimension zero, the SNC coordinate condition permits no member
through any point.  The genuine component family and the support of
$D$ are empty.  A coherent ideal with empty zero locus is the unit ideal.
Choose the unique function on the empty component family, $U$ equal to
the ambient manifold, and the empty product.  All assertions follow,
including the vacuous original-point coefficient statement.
\end{proof}

\begin{corollary}[The actual local divisor sum]
\label{hir:locally-finite-actual-divisor-sum-stalks}
\uses{hir:locally-finite-original-component-product,hir:finite-component-sum-support-containment}
\lean{Hironaka.HironakaLocallyFiniteDivisorSum.exists_locallyFinite_original_component_sums}
\leanok
With the same genuine global function $a$, each local list $A$ defines
an actual effective divisor $G_U=\sum_{(C,b)\in A}bC$.  At every point
$y\in U$, the complete stalks of $I_D$ and $I_{G_U}$ are equal.
Thus these sums give the locally finite effective-divisor dictionary
with the original divisor's actual multiplicities, without compactness.
\end{corollary}
\begin{proof}
Each retained carrier is an actual closed smooth hypersurface, so it
defines a genuine invertible hypersurface ideal and effective divisor.
The genuine finite-sum ideal identity is
$I_{G_U}=\prod_{(C,b)\in A}I_C^b$.  Substitute this identity in the
whole-stalk equivalence of the theorem.  All coefficients and carrier
conditions remain those produced there.  In dimension zero the local
list is empty and the actual sum is the zero divisor.
\end{proof}
